\documentclass[12pt]{article}

%Dois dielectricos na vertical
%Faro, Wed Mar 19 19:47:16 WET 2014
%Written by Tordar

\usepackage{graphicx}
\usepackage{pst-all}



\begin{document}

\pagestyle{empty}

\begin{pspicture}(-7,-7)(7,7)
\psset{xunit=1cm,yunit=1cm}
%Fios superiores:
\psline[linewidth=1mm,linecolor=black,linestyle=solid](0,4)(0,5)
\psline[linewidth=1mm,linecolor=black,linestyle=solid](-2.5,3)(-2.5,4)(2.5,4)(2.5,3)
%Fios inferiores:
\psline[linewidth=1mm,linecolor=black,linestyle=solid](0,-5)(0,-4)
\psline[linewidth=1mm,linecolor=black,linestyle=solid](-2.5,-3)(-2.5,-4)(2.5,-4)(2.5,-3)
%Dielectrico a esquerda:
\psframe[linewidth=1mm,linecolor=black,linestyle=solid](-4,-3)(-1,3)
%Dielectrico a direita:
\psframe[linewidth=1mm,linecolor=black,linestyle=solid](1,-3)(4,3)
%Seta dupla:
\psline[linewidth=1mm,linecolor=blue,linestyle=solid,arrows=<->](4.5,-3)(4.5,3)
\rput{0}(-3.5,3.5){\scalebox{1.5}{$S/2$}}
\rput{0}(3.5,3.5){\scalebox{1.5}{$S/2$}}
\rput{0}(5,0){\scalebox{1.5}{$d$}}
\rput{0}(-2.5,0){\scalebox{1.5}{$\varepsilon_1$}}
\rput{0}(2.5,0){\scalebox{1.5}{$\varepsilon_2$}}
\end{pspicture}

\end{document}
